% ATTEMPT_STATUS: unresolved
% TOP_PROBLEM: {"problemId": "problem.vassiliev-finite-type-invariant-separation-conjecture", "canonicalTitle": "Vassiliev Separation by Abelian-Valued Finite-Type Invariants", "releaseRank": 288, "primaryDomain": "geometry_topology", "primaryDomainLabel": "Geometry & topology", "displayStatus": "open", "releaseStatus": "open"}
% !TeX program = lualatex
\documentclass[11pt]{article}
\usepackage[margin=25mm]{geometry}
\usepackage{fontspec}
\setmainfont{Latin Modern Roman}
\usepackage{amsmath,amssymb,mathtools}
\usepackage{unicode-math}
\setmathfont{Latin Modern Math}
\usepackage{xurl,hyperref}
\hypersetup{colorlinks=true,urlcolor=blue}
\setcounter{secnumdepth}{0}
\setlength{\parindent}{0pt}
\setlength{\parskip}{5pt}
\setlength{\emergencystretch}{3em}
\begin{document}
\section*{\texorpdfstring{Vassiliev Separation by Abelian-Valued Finite-Type Invariants}{Vassiliev Separation by Abelian-Valued Finite-Type Invariants}}\label{288}
\subsection{Definitions and mathematical statement}
Let $\mathcal K$ be the set of smooth closed oriented unframed knots in $S^3$ modulo orientation-preserving ambient isotopy. An invariant $v:\mathcal K\to A$, with $A$ abelian, extends linearly to ${\mathbb{Z}}[\mathcal K]$ and to singular knots using
$v(K_\times)=v(K_+)-v(K_-)$. It has finite type at most $d$ if it vanishes on every singular knot with $d+1$ transverse double points. The target is
\[
 K_1\not\cong K_2\quad\Longrightarrow\quad
 \exists A,d,v\text{ of this type with }v(K_1)\ne v(K_2).
\]
The coefficient group and order may depend on the pair. A knot and its reverse are included if they are distinct oriented types. Rational-valued separation is not silently substituted for the arbitrary-abelian-coefficient version.
\par\smallskip\textit{Formulation and concept references:} \href{https://doi.org/10.2140/agt.2026.26.1037}{[S1]}.

\subsection{Short English statement}
Can any two distinct oriented knot types be distinguished by some finite-type invariant, allowing its finite order and abelian coefficient group to depend on the pair?
\subsection{Research attempt}
\paragraph{Scope and source check.}
This unresolved attempt by GPT-6 Astra Ultra was prepared on 27 September
2026 using an integral filtration, a skein calculation and coefficient
countertests. The coefficient group in the statement is essential. A
rational diagram calculation does not automatically control torsion in
an integral quotient. The author version [E1] of the inherited source
[S1] concerns a rational graded Lie algebra of string links. Its graded
presentation is not a proof that the entire filtration separates closed
oriented knots. Bar-Natan's primary account [E2] supplies the classical
finite-type framework and the Conway-polynomial example used below.

\paragraph{An exact universal reformulation with arbitrary coefficients.}
Let $V=\mathbb Z[\mathcal K]$, the free abelian group on the oriented
unframed knot types in the statement. For $r\geq1$, let $F_rV$ be the
subgroup generated by the alternating sums of the $2^r$ resolutions of
all singular knots with $r$ transverse double points. Set $F_0V=V$.
Resolving one selected double point first expresses a generator of
$F_{r+1}V$ as the difference of two generators of $F_rV$, so
\[
                   F_{r+1}V\subseteq F_rV.
\]
No rationalization or choice of chord-diagram basis enters this
inclusion. For each $d\geq0$ define
\[
       U_d=V/F_{d+1}V,\qquad v_d(K)=[K]\in U_d.
\]
This is an abelian-valued invariant of type at most $d$: its extended
value on every $(d+1)$-singular knot is zero by the definition of the
quotient. Conversely, any $A$-valued invariant $v$ of type at most $d$
extends to a homomorphism $V\to A$ killing $F_{d+1}V$. It therefore
factors uniquely through $U_d$.

Consequently, for a particular pair $K,L$ the existence of a separating
finite-type invariant is exactly
\[
       [K]-[L]\notin\bigcap_{d\geq0}F_{d+1}V.                 \tag{1}
\]
If the difference survives one quotient, that quotient itself supplies
the permitted target group. If it belongs to the intersection, every
finite-type invariant to every abelian group kills it. Equivalently,
the map
\[
          \mathcal K\longrightarrow\varprojlim_d U_d
\]
must be injective. The inverse limit here is an algebraic completion;
no analytic convergence of a power series has been established or is
needed for this reformulation.

There is a further distinction: injectivity of the linear map
$V\to\varprojlim_dU_d$ is sufficient but stronger than injectivity on
the knot basis. Distinct images of basis elements need not be linearly
independent. For example the homomorphism from the free abelian group on three
generators $e_0,e_1,e_2$ to $\mathbb Z$, sending them to $0,1,2$,
distinguishes these three basis elements while killing
$e_0-2e_1+e_2$. This is an algebraic warning, not a proposed knot
counterexample. A successful proof of (1) need only rule out nonzero
\emph{pair differences} in the intersection.

\paragraph{Why rational coefficients are a genuinely stronger request.}
Write $x_d$ for the image of $[K]-[L]$ in $U_d$. A homomorphism
$U_d\to\mathbb Q$ detects $x_d$ precisely when
$x_d\otimes1\ne0$ in $U_d\otimes\mathbb Q$. Indeed, one direction is
immediate; in the other, a linear functional on the rational vector
space can be chosen nonzero on the given vector. For an abelian group,
the kernel of its rationalization consists exactly of its torsion
elements. Thus a nonzero torsion $x_d$ is detectable with an allowed
abelian target but not with a rational target at that same order.

The difference can persist through an entire tower of abstract
quotients. Consider the filtration $G_d=2^d\mathbb Z$ on $\mathbb Z$.
Its intersection is zero, and reduction modulo $2^d$ eventually detects
every nonzero integer. Nevertheless
\[
              (\mathbb Z/2^d\mathbb Z)\otimes\mathbb Q=0
\]
for every $d$. All rational invariants factoring through these
quotients vanish. This filtration is not asserted to occur for knots;
it proves that separatedness before rationalization cannot be deduced
from a claim about rational associated graded groups alone. Conversely,
a proof of rational separation would imply the requested abelian
separation, since $\mathbb Q$ is an allowed target. The implication does
not justify replacing the target in advance.

\paragraph{A concrete separating calculation and its exact limit.}
The normalized Conway polynomial of an oriented link satisfies
\[
       \nabla_{K_+}(z)-\nabla_{K_-}(z)=z\nabla_{K_0}(z),
       \qquad \nabla_{\text{unknot}}=1,
\]
where $K_0$ is the oriented smoothing. Apply this identity at $r$
distinct singular points. The alternating sum of the Conway
polynomials of their resolutions is $z^r$ times the Conway polynomial
of the completely smoothed oriented link. Therefore
\[
                  a_d(K)=[z^d]\nabla_K(z)
\]
vanishes on every singular knot with $d+1$ double points. This proves
directly that $a_d$ is an integer-valued finite-type invariant of order
at most $d$. The argument is the standard skein proof in [E2]; it makes
no completeness claim for this family.

For a trefoil with Seifert matrix
\[
                  B=\begin{pmatrix}1&0\\-1&1\end{pmatrix},
\]
the normalized Alexander determinant is
\[
 \det(t^{1/2}B-t^{-1/2}B^T)
       =t+t^{-1}-1=1+(t^{1/2}-t^{-1/2})^2.
\]
With $z=t^{1/2}-t^{-1/2}$ this gives $\nabla=1+z^2$.
Thus $a_2$ distinguishes this trefoil from the unknot, with values $1$
and $0$. This is an actual instance of (1), not just an abstract
existence assertion about a quotient.

The same family has a sharp limitation. The mirror trefoil has Seifert
matrix $-B^T$ and the same determinant in this size, hence the same
Conway polynomial. The symmetric forms $B+B^T$ and $-(B+B^T)$ have
signatures $2$ and $-2$, respectively: the first has eigenvalues $1,3$.
Using the standard invariance of the knot signature, the two knots are
distinct. Consequently every Conway coefficient agrees on this pair
although a different knot invariant separates them. This does not say
that all finite-type invariants agree; it only disproves completeness
of this particular finite-type family. Reversing the orientation of a
knot also preserves its Conway polynomial. The root statement includes
oriented reverse pairs, so an argument restricted to such polynomial
data would have a second unaddressed limitation.

\paragraph{Why a graded calculation does not finish the filtration argument.}
The groups $F_dV/F_{d+1}V$ describe successive layers, and relations
among diagrams can identify those layers or parts of them. But knowing
all successive quotients does not by itself prove that the filtration
has zero intersection. As an elementary example, replace any separated
filtered group $W$ by $W\oplus C$ and put
$\widetilde F_d=F_dW\oplus C$ for every $d$, with $C\ne0$. The
successive quotients are unchanged, while the new intersection contains
$C$. This does not construct a hidden subgroup in the knot filtration.
It specifies the logical information absent from a graded
presentation. A reconstruction theorem must separately exclude such
invisible elements, at least among the differences of knot types.

Likewise, pairwise separation is not equivalent to a single order
working for every pair. In the elementary tower
$\mathbb Z\to\mathbb Z/2^d\mathbb Z$, all orders together separate
integers but no one order does. For knots the allowed order may depend
on $K,L$. A computation through order $D$ can certify separation if one
computed invariant differs. Agreement through order $D$ is consistent
both with a first difference at a higher order and with agreement at
every order. Neither alternative follows from a finite table.

This also identifies the needed role of a prospective complexity
bound. Suppose one had a proved function $b(c)$ such that every two
nonisotopic knots with diagrams of at most $c$ crossings were separated
by order at most $b(c)$, together with an effective exact description
of the relevant integral quotients. Then a finite computation could
settle separation for each bounded diagram class. The bound is a new
substantive hypothesis, not a consequence of the finite number of
those diagrams: finiteness allows taking a maximum only after one
already knows that each distinct pair is separated. Using that maximum
to prove separation would be circular.

\paragraph{Verification, first gap and outcome.}
The local checks expand the displayed determinant, verify the two
signature matrices, and test the alternating finite-difference
cancellation underlying finite-type order. The filtration inclusions,
universal property and rationalization distinction are proved above;
finite arithmetic does not certify the infinite intersection in (1).
The source review preserves the difference between closed oriented
knots and the rational graded string-link results in [S1,E1]. No
claimed solution of the separation problem is used as an input.

The first unresolved step is now precise: exclude $[K]-[L]$ from the
intersection of the integral Vassiliev filtration whenever $K,L$ are
distinct oriented types, or exhibit such a pair in the intersection.
The trefoil calculation proves a restricted example, and the mirror
calculation prevents promoting that example to a universal method.
\textbf{Outcome: UNRESOLVED.} The attempt gives an exact reformulation
for the requested coefficient scope, not a proof of separation.
\subsection{Sources}
\begin{itemize}
\item[S1]Bruno Dular, Primitive Feynman Diagrams and the Rational Goussarov–Habiro Lie Algebra of String Links, AGT 26:3 (2026), pp. 1037–1038, definitions and m=1 separation question; Ohtsuki problem list for oriented-knot formulation.
\url{https://doi.org/10.2140/agt.2026.26.1037}.
\textit{Formulation source.}
\item[E1]Bruno Dular, \emph{Primitive Feynman diagrams and the rational Goussarov--Habiro Lie algebra of string links}, author version of [S1], arXiv:2406.01093.
\url{https://arxiv.org/abs/2406.01093}.
\textit{Consulted 27 September 2026 for the distinction between a rational graded presentation and full knot separation.}
\item[E2]Dror Bar-Natan, \emph{On the Vassiliev knot invariants}, Topology 34 (1995), 423--472, especially Problem 1.1 and Section 2.1.1.
\url{https://www.math.toronto.edu/~drorbn/papers/OnVassiliev/OnVassiliev.pdf}.
\textit{Primary author text for the finite-type framework and the Conway skein calculation; consulted 27 September 2026.}
\end{itemize}
\end{document}
